\section{Method: The VIGIL Framework}
\label{sec:method}

\subsection{Overview}
\label{sec:overview}

\sysname{} comprises two layers.
A sensing pipeline discovers AI-litigation filings in CourtListener/RECAP~\citep{courtlistener2023} and legal news and converts them into case records; this engineering infrastructure is documented in the supplementary material.
The remainder of this section describes the evaluated NLP components: LegalRisk-RAG (\S\ref{sec:rag}), the goal-conditioned agentic loop (\S\ref{sec:agent}), the temporal precedent graph (\S\ref{sec:temporal}), and the \plre{} decision procedure (\S\ref{sec:plre}).

\subsection{LegalRisk-RAG}
\label{sec:rag}

The LegalRisk-RAG pipeline adapts standard RAG to legal text in three respects.
\textbf{Chunking} segments documents at logical structural boundaries---numbered claims, ``WHEREFORE'' clauses, section headings---rather than fixed token windows; each chunk retains a metadata header (case identifier, filing date, court level, section type).
\textbf{Embedding:} we compared a dense retriever (text-embedding-3-large), RECAP-tuned LegalBERT~\citep{chalkidis2020legalbert}, and a hybrid BM25~\citep{robertson2009bm25} + neural re-ranker on a 200-pair development set; the hybrid performed best and is the default.
\textbf{Dynamic context expansion (precedent chaining):} when a retrieved chunk references a prior case, the pipeline appends that opinion as secondary context and queries the temporal precedent graph for subsequent treatment, so that the Analyzer receives both semantic neighbors and validity metadata.

\subsection{Goal-Conditioned Agentic Loop}
\label{sec:agent}

The four specialized agents and the central Goal Checker operate as described in the supplementary material.
The Planner decomposes unmet sub-goals into ordered retrieval/analysis sub-tasks; the Retriever runs LegalRisk-RAG queries and reformulates via legal entity term expansion when $\mathcal{G}_{\mathrm{evidence}}$ is unmet; the Analyzer computes $\mathrm{RiskScore}(d_i, P, C)$, clusters evidence, and maps clusters to a risk taxonomy (Training Data, Output Reproduction, API Misuse, Model Weights); the Goal Checker evaluates $\mathcal{G}(R_t)$ after each pass and returns machine-readable residuals; the Planner converts these critiques into expanded entity/date/theory queries; and the Reporter assembles the risk register, per-risk evidence with citations, and mitigation guidance.

\subsection{Temporal Precedent Graph}
\label{sec:temporal}

\sysname{}'s central component is a treatment layer that determines when each authority ceases to be operative; a simple retrieval-time gate then applies it.
The layer is realized as a temporal precedent graph $P_t = (V, E_t)$.
Each node $v \in V$ denotes a case, order, or doctrinal rule and stores jurisdiction, court level, issue tags, filing and decision dates, and a validity interval $[s_v, e_v)$, where $e_v$ remains open unless the node is later vacated, overruled, or superseded; material narrowing instead lowers its treatment weight.
Directed edges encode treatment relations: \textsc{Cites}, \textsc{Distinguishes}, \textsc{Narrows}, \textsc{Supersedes}, and \textsc{Overrules}.
Retrieval for a query at date $t_q$ first obtains semantically relevant candidates and then applies validity gating:
\begin{equation}
\label{eq:gate}
\begin{aligned}
\mathrm{score}(v, q, t_q) ={}& \mathrm{sim}(v, q) \cdot \mathbf{1}[s_v \le t_q < e_v] \\
&\cdot w(\mathrm{treat}(v, t_q)).
\end{aligned}
\end{equation}
Here $w(\cdot)$ is a fixed monotone schedule: positively treated or still-open authorities receive weight 1.0; cautionary authorities receive 0.75; material \textsc{Narrows} receives 0.50; fully closed (overruled or superseded) authorities are already excluded by the indicator.
The operator separates legal authority from topical similarity: a stale but lexically close case may still appear as context, yet it cannot support a risk decision without passing the validity check.
Equation~\eqref{eq:gate} is therefore a constrained retrieval operator over $(q, t_q)$: a host retriever scores semantic relevance as usual, and \plre{} receives only support whose legal status remains operative at $t_q$.

\paragraph{Why not a date filter?}
Once intervals are known, Eq.~\eqref{eq:gate} is a metadata filter any vector store can apply.
The difficulty is upstream and conceptual: $e_v$ is set by later documents that overrule, supersede, or narrow $v$, so it is not recoverable from $v$'s own filing date.
A pure recency or document-date filter therefore cannot implement the same operator; the empirical question is whether treatment edges can be extracted reliably enough to use.
Table~\ref{tab:treatment} addresses that question: negative-treatment micro-F1 is 0.80 on 600 gold spans and 0.81 against a commercial citator (Protocol~P2; missed-action rate 20.4\%; 3.8\% of held-out \plre{} pairs flip after citator correction).

\begin{table}[t]
\centering
\small
\begin{adjustbox}{max width=\columnwidth}
\begin{tabular}{lccc}
\toprule
Reference & P & R & F1 \\
\midrule
Gold spans ($n = 600$) & 0.84 & 0.77 & 0.80 \\
Citator (Protocol P2) & 0.83 & 0.80 & 0.81 \\
\bottomrule
\end{tabular}
\end{adjustbox}
\caption{Treatment-layer validation (neg.-treatment micro over \textsc{Overrules}/\textsc{Supersedes}/material \textsc{Narrows}).}
\label{tab:treatment}
\end{table}

\paragraph{Graph construction.}
Nodes are induced from CourtListener/RECAP with intervals initialized to $[s_v, \infty)$.
Edges are extracted by a citation-span detector plus an LLM classifier (GPT-4o, $T = 0$); interval-closing edges must be anchored to a detected span, decision date, and source document, and gold evaluation uses spans held out from prompts.
A confirmed \textsc{Overrules}/\textsc{Supersedes} edge at $t$ sets $e_v \leftarrow t$; material \textsc{Narrows} sets $w = 0.50$ from $t$.
The graph has 1{,}812 nodes and 4{,}760 edges.
Low-confidence negative-treatment edges are routed to human review rather than auto-closing intervals.

\subsection{PLRE Decision Procedure}
\label{sec:plre}

We cast product-litigation alignment as a risk-entailment task.
Given a product spec $P$ with attributes $\{p_1, \ldots, p_m\}$, a legal-risk theory $T_j$, and a query date $t_q$, \plre{} predicts
\begin{equation}
y_{ij} = \mathbf{1}[P, p_i, P_{t_q} \models T_j].
\end{equation}
Unlike symmetric similarity, \plre{} is directional.
A pair $(p_i, T_j)$ is promoted to the risk register only when (i)~the entailment score exceeds the validation-set threshold, (ii)~the top supporting precedents remain valid at $t_q$, and (iii)~the report includes mitigation text grounded in at least three cited authorities when available.

\plre{} gold labels are constructed from anonymized product-attribute cards and docket-grounded theory cards.
Annotators judge a directional pair as positive only when the product attribute, legal theory, query date, and at least one valid supporting authority jointly establish exposure; negative pairs comprise ordinary mismatches and stale-authority distractors.
Critically, the validity judgment encoded in the labels is produced by annotators reading docket and opinion text, independently of \sysname{}'s automatically extracted treatment graph, so the gate must recover label validity rather than presuppose it.
Product archetypes are tracked in the split manifest, and year and litigation type are stratified to reduce temporal and topical leakage.
The formulation is system-agnostic: \plre{} is directional entailment subject to an exogenous, time-indexed support constraint.
